%%%% LISTA DE ABREVIATURAS E SIGLAS 
%%
%% Relação, em ordem alfabética, das abreviaturas (representação de uma palavra por meio de alguma(s) de sua(s) sílaba(s) ou
%% letra(s)), siglas (conjunto de letras iniciais dos vocábulos e/ou números que representa um determinado nome) e acrônimos
%% (conjunto de letras iniciais dos vocábulos e/ou números que representa um determinado nome, formando uma palavra pronunciável).
%%
%% Este arquivo para definição de abreviaturas, siglas e acrônimos é utilizado com a opção "glossaries" (pacote)

%% Abreviaturas: \abreviatura{rótulo}{representação}{definição}

\abreviatura{art.}{art.}{Artigo}
\abreviatura{cap.}{cap.}{Capítulo}
\abreviatura{sec.}{sec.}{Seção}

%% Siglas: \sigla{rótulo}{representação}{definição}

\sigla{cve}{CVE}{Vulnerabilidades e Exposições Comuns, do inglês \textit{Common Vulnerabilities and Exposures}}
\sigla{ioc}{IOC}{Indicador de Comprometimento, do inglês \textit{Indicator of Compromise}}
\sigla{siem}{SIEM}{Sistema de Gerenciamento de Informações e Eventos de Segurança, do inglês \textit{Security Information and Event Management}}
\sigla{llm}{LLM}{Modelo de linguagem de grande escala, do inglês, \textit{Large Language Model}}
\sigla{ips}{IPS}{\textit{Intrusion Prevention System}}
\sigla{ia}{IA}{Inteligência Artificial, do inglês \textit{Artificial Intelligence}}
\sigla{ml}{ML}{Aprendizado de Máquina, do inglês \textit{Machine Learning}}
\sigla{nlp}{NLP}{Processamento de Linguagem Natural, do inglês \textit{Natural Language Processing}}
\sigla{soc}{SOC}{Centro de Operações de Segurança, do inglês \textit{Security Operations Center}}
\sigla{gdpr}{GDPR}{Regulamento Geral de Proteção de Dados, do inglês \textit{General Data Protection Regulation}}
\sigla{pci-dss}{PCI-DSS}{Padrão de Segurança de Dados da Indústria de Cartões de Pagamento, do inglês \textit{Payment Card Industry Data Security Standard}}
\sigla{cwe}{CWE}{Common Weakness Enumeration, do inglês \textit{Common Weakness Enumeration}}
\sigla{nist}{NIST}{Instituto Nacional de Padrões e Tecnologia, do inglês \textit{National Institute of Standards and Technology}}
\sigla{cisa}{CISA}{Agência de Segurança de Infraestrutura e Cibernética, do inglês \textit{Cybersecurity and Infrastructure Security Agency}}
\sigla{nvd}{NVD}{Banco Nacional de Vulnerabilidades, do inglês \textit{National Vulnerability Database}}
\sigla{cvss}{CVSS}{Sistema de Pontuação de Vulnerabilidades Comuns, do inglês \textit{Common Vulnerability Scoring System}}
\sigla{pln}{PLN}{Processamento de Linguagem Natural, do inglês \textit{Natural Language Processing}}
\sigla{ri}{RI}{Resposta a Incidentes, do inglês \textit{Incident Response}}
\sigla{ei}{EI}{Inteligência de Eventos, do inglês \textit{Event Intelligence}}
\sigla{ta}{TA}{Análise de Ameaças, do inglês \textit{Threat Analysis}}
\sigla{dvwa}{DVWA}{Damn Vulnerable Web Application, do inglês \textit{Damn Vulnerable Web Application}}
\sigla{cnn}{CNN}{Rede Neural Convolucional, do inglês \textit{Convolutional Neural Network}}
\sigla{svm}{SVM}{Máquina de Vetores de Suporte, do inglês \textit{Support Vector Machine}}
\sigla{ip}{IP}{Protocolo de Internet, do inglês \textit{Internet Protocol}}
\sigla{tf-idf}{TF-IDF}{Frequência de Termo–Frequência Inversa de Documento, do inglês \textit{Term Frequency–Inverse Document Frequency}}

\sigla{oscti}{OSCTI}{\textit{Open-Source Cyber Threat Intelligence}}
\sigla{spl}{SPL}{\textit{Splunk Processing Language}}
\sigla{kql}{KQL}{\textit{Kusto Query Language}}
\sigla{cot}{CoT}{\textit{Chain-of-Thought}}
\sigla{cotr}{CoT-R}{\textit{Chain-of-Thought + Reflection}}
\sigla{ir}{IR}{\textit{Intermediate Representation}}
\sigla{rag}{RAG}{\textit{Retrieval-Augmented Generation}, em português, Geração Aumentada por Recuperação}
\sigla{bgl}{BGL}{\textit{BlueGene/L}}
\sigla{dvwa}{DVWA}{\textit{Damn Vulnerable Web Application}}
\sigla{gpu}{GPU}{\textit{Graphics Processing Unit}}
\sigla{ids}{IDS}{\textit{Intrusion Detection System}}
\sigla{epp}{EPP}{\textit{Endpoint Protection Platforms}}
\sigla{pci-dss}{PCI-DSS}{\textit{Payment Card Industry Data Security Standard}}
\sigla{ueba}{UEBA}{\textit{User and Entity Behavior Analytics}}


%% LEIA:

%% Para usar o \gls, você deve colocar a sigla aqui em \sigla

%% ADICIONAR SIGLAS: Quando você inclui alguma sigla, pode ser necessário compilar umas duas vezes para essa aparecer na lista de siglas.

%% ATENÇÃO REMOVER SIGLAS: se você remover a sigla do seu texto (não for usar mais), você deve comentar essa aqui e remover os \gls{} dessa sigla (se não vai ficar aparecendo a sigla na lista). Em caso de ERRO, quando o LaTeX informa que você ainda tem a sigla no texto, mesmo que não tenha. Você deve limpar o cache - no OverLeaf, clique no erro, vá:
%  ->view error
%%   ->(role para baixo, até o final)
%%     ->e clique em Clear cached files


%% Acrônimos: \acronimo{rótulo}{representação}{definição}
%\acronimo{gimp}{Gimp}{Programa de Manipulação de Imagem GNU, do inglês \textit{GNU Image Manipulation Program}}
